\section{The Cloud Computing Parallel, 2006--2020}\label{sec:cloud}

% Role: empirical analogue that disciplines the entry-cost channel (kappa_e, kappa_f in sec:ai-model).

\subsection{Why cloud computing is the right analogue}\label{sec:cloud-why}

Before commercial cloud services, a new technology firm had to buy and operate its own servers before it could test whether its product had demand.
The public launch of Amazon Web Services' storage and compute services in 2006 allowed firms to rent computing capacity by the unit instead, converting a fixed, sunk investment into a variable cost.
\citet{ewens2018cost} exploit this change and show that lower costs of experimentation altered how venture investors financed early-stage firms in the affected sectors.
In the model of \cref{sec:theory}, this shock maps naturally into a reduction of the entry cost $c_e$ and the fixed operating cost $c_f$, with no direct change in the labor requirement of incumbents.
This is the sense in which cloud computing isolates one of the two channels of the AI shock.
\placeholder{Add evidence on cloud adoption by firm size and age, for example \citet{jin2017economies}, and on the timing of diffusion in Europe relative to the United States.}

\subsection{Data}\label{sec:cloud-data}

We draw on multiple sources to characterize the cloud era.
For the United States: the Census Bureau's Business Dynamics Statistics (BDS) provides firm and establishment entry, exit, and employment by firm size and age annually from 1978--2021 \citep{bds2021}; the Business Formation Statistics (BFS) track business applications and planned-wage intentions from 2004 \citep{bfs2026}.
For Europe: Eurostat Structural Business Statistics and business demography statistics cover EU member states, and the OECD DynEmp database provides harmonized firm-level aggregates.

Size-class definitions differ across sources.
The upper bound of the band $\Bc$ aligns with the firm-size classes in the US data (5,000 employees), whereas European aggregate statistics typically group all firms with 250 or more employees, so the European analysis uses a narrower band where necessary.

\subsection{Empirical findings: cloud produced concentration, not offset}\label{sec:cloud-findings}

The 2006--2020 cloud era provides a cautionary baseline.
Contrary to the hypothesis that lower fixed costs would proliferate small firms, US employment became more concentrated \citep{bds2021}:

\begin{itemize}
  \item \textbf{Small-firm employment share fell}: from 50.9\% (2006) to 47.0\% (2019) for firms with fewer than 500 employees, and about 90\% of the net rise in employment over the period occurred at firms with 500 or more employees.
  \item \textbf{The startup rate declined}: the employer startup rate fell from 10.3\% (2006) to 8.5\% (2019), and the employment share of young firms (five years or younger) fell from 13.4\% to 10.6\%, continuing a secular decline that predates cloud computing \citep{decker2016declining}.
  \item \textbf{The largest firms gained share}: firms with 10,000 or more employees raised their share of employment from 26.3\% to 29.2\%, consistent with the rise of superstar firms documented by \citet{autor2020fall}.
  \item \textbf{Formation shifted to the nonemployer margin}: business applications rose 34\% between 2006 and 2019, while applications with planned wages fell 48\% \citep{bfs2026}.
\end{itemize}

The most cloud-exposed sector is a partial exception.
The startup rate in Information (NAICS 51) recovered from 8.5\% in 2010 to 11.3\% in 2019, above its 2006 level of 10.6\%, whereas professional services, retail, and manufacturing remained below their 2006 levels \citep{bds2021}.
Cloud computing also enabled notable successes (Airbnb, Stripe, Dropbox), but these firms ``graduated'' into the large-firm category rather than remaining in the small-firm band.

\subsection{Why did cloud fail to offset?}\label{sec:cloud-mechanism}

Our primary use of the cloud era is as a calibration target rather than a causal identification exercise.
We ask what parameter configuration of the model replicates the cloud-era pattern of lower computing costs combined with rising concentration and falling dynamism.
No single entry-side parameter does so: varying the entry cost, the fixed cost, entrant quality, or the founder pool on its own leaves the small-firm share rising.
Matching all three moments requires two ingredients together:

\begin{enumerate}
  \item \textbf{Labor-complementary technology} ($\sigma < \theta \approx 2.78$): adopting firms expand rather than shed labor, so employment shifts toward large adopters. This ingredient alone produces concentration but \emph{raises} the startup rate.
  \item \textbf{Entry headwinds}: a higher net employer entry cost, a smaller pool of potential entrants, and no productivity gain for cloud-native entrants, consistent with the regulatory, demographic, and labor-market forces behind declining dynamism \citep{akcigit2021ten}.
\end{enumerate}

\Cref{tab:cloud-results} reports the best-fitting configuration, which sets $\sigma = 2.0$, raises the employer entry cost by 20\%, and shrinks the founder pool by 20\%.
The two ingredients separate cleanly: the technology parameter drives the size distribution, while the entry-side headwinds drive the decline in dynamism.
Because the model's pre-shock levels are not calibrated to these moments, we compare changes rather than levels.

\begin{table}[t]
  \centering
  \caption{Best-fitting cloud-era configuration: model versus data, 2006--2019}
  \label{tab:cloud-results}
  \begin{tabular}{lcc}
    \toprule
    \multicolumn{3}{l}{\textit{Parameters}} \\
    $\sigma$ (substitution) & \multicolumn{2}{c}{$2.0$ (labor-complementary)} \\
    Employer entry cost & \multicolumn{2}{c}{$+20\%$ (headwinds offset savings)} \\
    Entrant productivity shift & \multicolumn{2}{c}{$0$} \\
    Founder pool & \multicolumn{2}{c}{$-20\%$ (declining dynamism)} \\
    \midrule
    Moment (change) & Model & Data \\
    \midrule
    Employment share of firms $<500$ (pp) & $-2.3$ & $-3.9$ \\
    Startup rate (pp) & $-1.6$ & $-1.8$ \\
    Young-firm employment share (pp) & $-3.1$ & $-2.8$ \\
    \bottomrule
  \end{tabular}
  \smallskip
  \footnotesize\textit{Notes: Model values average years 11--15 after the shock (centered on 2019) across 24 seeds, relative to a no-shock counterfactual. Data from \citet{bds2021}.}
\end{table}

The implication for AI is twofold.
First, even labor-complementary AI ($\sigma_{\text{ai}} < 2.78$) could produce concentration if entry headwinds persist.
Second, the offset found in \cref{sec:simulation} rests on an intact entry pipeline: imposing cloud-era founder headwinds on the baseline AI shock ($\sigma_{\text{ai}} = 3.5$, a 20\% cut in entry costs) turns a long-run net employment gain of about $+4\%$ into a loss of $-8.7\%$ after 13 years and $-15.4\%$ in the long run.

\subsection{Empirical strategy}\label{sec:cloud-strategy}

The calibration exercise disciplines the entry-cost channel while remaining agnostic about identification.
For readers interested in the causal effect of cloud computing on entry, the canonical event-study design compares industries by their ex ante exposure,
\begin{equation}\label{eq:cloud-event}
  y_{jt} = \alpha_j + \lambda_t + \sum_{k \neq -1} \beta_k \, \bigl( \text{Exposure}_j \times \mathbf{1}\{t = 2006 + k\} \bigr) + X_{jt}'\gamma + \varepsilon_{jt},
\end{equation}
where $j$ indexes industries, $y_{jt}$ is an outcome such as the firm entry rate or the employment share of the band $\Bc$, and $\text{Exposure}_j$ measures pre-2006 reliance on computing infrastructure (e.g., IT capital intensity from the BEA).
Because the aggregate US startup rate declined over this period \citep{decker2016declining}, identification comes from differences across industries relative to a common trend, not from changes in the aggregate series.
